\subsection{Hopf spaces}

DEFINITION 2. --- A Hopf space is a pointed topological space (X, e) equipped with a continuous composition law m: X × X → X such that

\begin{enumerate}
\def\labelenumi{(\roman{enumi})}
\item
  \(m(e, e) = e\) ;
\item
  there exist pointed homotopies at \(e\) relating the maps \(x \mapsto m(x, e)\) and \(x \mapsto m(e, x)\) to the map \(\mathrm{Id}_{\mathrm{X}}\) .
\end{enumerate}

One sometimes expresses properties (i) and (ii) by saying that e is a neutral element up to homotopy for the composition law m.

Note that there may exist several identity elements up to homotopy for a continuous composition law m on a topological space X. For example, for \(X = \mathbf{R}\) and \(m(x, y) = (x + y)/2\) , every \(x \in \mathbf{R}\) is a neutral element up to homotopy.

Examples. --- 1) Let G be a topological group and m its composition law. The identity element e of G is an identity element up to homotopy, and (G, e) is a Hopf space.

\begin{enumerate}
\def\labelenumi{\arabic{enumi})}
\setcounter{enumi}{1}
\tightlist
\item
  Let X be a topological space and let x be a point of X. Equipped with the concatenation of loops at x, the pointed space \((\Omega_{x}(\mathrm{X}), e_{x})\) is a Hopf space. Indeed, one first has \(e_{x} * e_{x} = e_{x}\). On the other hand, let \(\psi: I \to I\) be the function defined by \(\psi(t) = 2t\) for \(0 \leqslant t \leqslant \frac{1}{2}\) and \(\psi(t) = 1\) for \(\frac{1}{2} \leqslant t \leqslant 1\) (cf. III, p.~291), and let \(\sigma: I \times I \to I\) be a strict homotopy joining \(\psi\) to \(\mathrm{Id}_{\mathrm{I}}\) (III, p.~289, example). Then, for every loop \(c \in \Omega_{x}(\mathrm{X})\) , the map \(c \circ \sigma: I \times I \to X\) is a strict homotopy joining \(c * e_{x}\) to c in \(\Omega_{x}(\mathrm{X})\) .
\end{enumerate}

Let \(\tau\) be the map \(\Omega_x(\mathrm{X}) \times \mathbf{I} \to \Omega_x(\mathrm{X})\) defined by \(\tau(c, s)(t) = c \circ \sigma(s, t)\) . Let us prove that \(\tau\) is continuous. By Proposition 1 of III, p.~257, it suffices to show that the map \((c, s, t) \mapsto c(\sigma(s, t))\) from \(\Omega_x(\mathrm{X}) \times \mathbf{I} \times \mathbf{I}\) into X is continuous, that is, since \(\mathbf{I} \times \mathbf{I}\) is compact, that the map \(c \mapsto c \circ \sigma\) from \(\mathcal{C}_c(\mathbf{I};\mathrm{X})\) into \(\mathcal{C}_c(\mathbf{I} \times \mathbf{I};\mathrm{X})\) is continuous. This last assertion then follows from the lemma, I, p.~132, b). Consequently, the map \(\tau\) is a homotopy connecting the map \(c \mapsto c * e_x\) to the identity map of \(\Omega_x(\mathrm{X})\) . We have \(\tau(e_x, s)(t) = x\) for all \(s, t \in I\) , therefore \(\tau\) is a pointed homotopy at \(e_x\). One argues similarly for the map \(c \mapsto e_x * c\) .

PROPOSITION 3. — Let (X, e) be a Hopf space and m: X × X → X its composition law. For all loops c, c′ in X at e, one has:

\[
c ^ {\prime} * c \sim c * c ^ {\prime} \sim m \circ (c, c ^ {\prime}),
\]

where ∼ denotes the strict homotopy relation. The composition law of the group \(\pi_{1}(\mathrm{X},e)\) is the composite of the canonical isomorphism \(\pi_{1}(\mathrm{X},e)\times\pi_{1}(\mathrm{X},e)\to\pi_{1}(\mathrm{X}\times\mathrm{X},(e,e))\) (III, p.~297, corollary) and the homomorphism \(\pi_{1}(m,(e,e))\) from \(\pi_{1}(\mathrm{X}\times\mathrm{X},(e,e))\) into \(\pi_{1}(\mathrm{X},e)\) . The group \(\pi_{1}(\mathrm{X},e)\) is commutative.

Let \(\mu\colon\pi_{1}(\mathrm{X},e)\times\pi_{1}(\mathrm{X},e)\to\pi_{1}(\mathrm{X},e)\) be the group homomorphism given by the composite of the canonical isomorphism of \(\pi_{1}(\mathrm{X},e)\times\pi_{1}(\mathrm{X},e)\) onto \(\pi_{1}(\mathrm{X}\times\mathrm{X},(e,e))\) and the homomorphism \(\pi_{1}(m,(e,e))\) . It remains to prove that one has

\[
\mu (\gamma , \gamma^ {\prime}) = \gamma \gamma^ {\prime} = \gamma^ {\prime} \gamma\tag{1}
\]

for all \(\gamma, \gamma' \in \pi_1(X, e)\) . In view of remark 2 of III, p.~297, we have:

\[
\mu (\gamma , \gamma^ {\prime}) = \mu (\gamma , \varepsilon_ {e}) \mu (\varepsilon_ {e}, \gamma^ {\prime}) = \mu (\varepsilon_ {e}, \gamma^ {\prime}) \mu (\gamma , \varepsilon_ {e}).\tag{2}
\]

Let \(c \in \Omega_e(\mathrm{X})\) be a loop of class \(\gamma\) , and denote by \(m_1 \colon \mathrm{X} \to \mathrm{X}\) the map defined by \(m_1(x) = m(x, e)\) ; then, \(\mu(\gamma, \varepsilon_e)\) is the class of \(m_1 \circ c\) . By definition of a Hopf space, the pointed maps at \(e\) , \(m_1\) and \(\mathrm{Id_X}\) have the same pointed homotopy class at \(e\) . Consequently, the loops at \(e\) , \(m_1 \circ c\) and \(c\) are strictly homotopic (III, p.~296, cor. 1 of prop. 3) and one has \(\mu(\gamma, \varepsilon_e) = \gamma\) . One proves in an analogous way that one has \(\mu(\varepsilon_e, \gamma') = \gamma'\) . Relation (1) then follows from relation (2).

Note. --- Let \((X,e)\) be a Hopf space. By prop. 3, the composition map of path classes, \(\pi_{1}(X,e)\times\pi_{1}(X,e)\to\pi_{1}(X,e)\) is continuous if \(\pi_{1}(X,e)\) is equipped with the quotient topology of the compact-convergence topology on \(\Lambda_{e}(X)\) . This is indeed the composite of the continuous isomorphism \(\pi_{1}(X,e)\times\pi_{1}(X,e)\to\pi_{1}(X\times X,(e,e))\) (III, p.~298, remark 3) and the continuous map \(m_{*}:\pi_{1}(X\times X,(e,e))\to\pi_{1}(X,e)\) (III, p.~294, remark 1).

